% Keep the TeX year register available for PGF; the template redefines \year.
\let\CalendarYear\year
\documentclass{ieeeaccess}
\let\year\CalendarYear
\usepackage{cite}
\usepackage{amsmath,amssymb}
\usepackage{graphicx}
\usepackage{booktabs,tabularx}
\usepackage{tikz}
\usetikzlibrary{arrows.meta,positioning}
\usepackage[hidelinks]{hyperref}

% The supplied IEEE Access class expects this dimension even for TikZ figures.
\makeatletter
\@ifundefined{xfigwd}{\newdimen\xfigwd}{}
% Preserve the template headers without implying that this draft is published.
\let\access@titlepage\ps@titlepage
\let\access@headings\ps@headings
\newcommand{\draft@footers}{%
  \def\@oddfoot{\normalfont\scriptsize MANUSCRIPT DRAFT\hfill\thepage}%
  \def\@evenfoot{\normalfont\scriptsize\thepage\hfill MANUSCRIPT DRAFT}}
\def\ps@titlepage{\access@titlepage\draft@footers}
\def\ps@headings{\access@headings\draft@footers}
\makeatother
\setlength{\xfigwd}{0pt}
% Reserve room inside the title block for the template's abstract bullets.
\newlength{\paperabstractwidth}
\setlength{\paperabstractwidth}{\dimexpr\titlewidth-10pt\relax}
\let\abstractwidth\paperabstractwidth

\begin{document}
\history{Manuscript draft.}
\doi{not assigned}

\title{FalconVQA: Timestamp-Grounded Search and Question Answering for Video Repositories}
\author{VISHAL K. BHAT\authorrefmark{1},
TALLAM SRI SAI SUBRAMANYAM\authorrefmark{1},
MEDHA SANKETH\authorrefmark{1}, \mbox{AND AVINASH ANISH\authorrefmark{2}}}
\address[1]{Department of Computer Science and Engineering, RV College of Engineering, Bengaluru 560059, India}
\address[2]{Department of Information Science and Engineering, RV College of Engineering, Bengaluru 560059, India}
\corresp{Correspondence: Vishal K. Bhat (e-mail: vishalkbhat.cs23@rvce.edu.in).}
\markboth{Bhat \MakeLowercase{\textit{et al.}}: FalconVQA}{Bhat \MakeLowercase{\textit{et al.}}: FalconVQA}

\begin{abstract}
Finding an event in a video repository requires locating the evidence in the recording, not merely identifying the file or generating a plausible description. We present FalconVQA, a system that indexes timestamped visual descriptions and speech transcripts, retrieves ranked video intervals, and presents evidence as playable clips. Its VIDRA pipeline produces two views of a recording: a moment index for local evidence and video-level aggregates for broader questions. A conversational interface retrieves from the user's project and presents answers alongside source intervals. In a component pilot on an Apple M2, scene-based segmentation matched all nine deliberately abrupt transitions across three synthetic videos within 250~ms and took 47.8~ms per video at the median. The fixed 20-s grid matched none of the deliberately off-grid transitions. Local retrieval over 100, 1,000, and 5,000 synthetic 1,536-dimensional vectors took 12.2, 144.4, and 770.9~ms at the median, respectively, excluding query embedding and network traffic. These measurements characterize two components, not semantic retrieval accuracy or end-to-end answer quality. We distinguish them from exploratory measurements of an earlier backend and specify how to evaluate the current system on real footage with human-annotated evidence.
\end{abstract}

\begin{keywords}
Multimodal retrieval, retrieval-augmented generation, temporal grounding, video moment retrieval, video question answering.
\end{keywords}

\titlepgskip=-15pt
\maketitle
\pagestyle{headings}

\section{Introduction}
\label{sec:intro}
\PARstart{A}{n} operator reviewing an archive may know that a delivery was discussed, or that a container was moved, without knowing the recording or the timestamp. Filename and date filters do not express such a request. A transcript can locate spoken content but cannot describe a silent action; a frame-based description can record visible objects but omit speech, timing, or the relationship between two actions. A search interface must therefore handle both modalities and make it possible to verify where an answer came from.

Video moment retrieval and video question answering address different parts of this problem. The first returns a time interval for a description, while the second returns a textual response that can require reasoning over several intervals \cite{qvhighlights,nextqa,egoschema}. In repository search, a query may also span several recordings, and a user may want to inspect the original footage before acting on an answer. We treat \emph{the playable interval}, not the generated text alone, as the primary unit of evidence.

FalconVQA combines an offline indexing pipeline with an interactive retrieval interface. Rather than processing raw video for every question, it aligns picture and speech to a shared timeline and indexes short representations of that evidence. The query service can retrieve moments directly or combine them with precomputed summaries and chapters. A separate application layer translates conversational requests into scoped tool calls and renders the selected intervals.

This work makes three contributions:
\begin{enumerate}
\item An implemented video indexing pipeline with a shared temporal representation for visual descriptions, speech, and video-level aggregates.
\item A repository workflow in which retrieval results preserve recording identity and playable time bounds, and video tools resolve evidence within the current project.
\item A reproducible component pilot that measures the costs of segmentation and file-backed retrieval, plus an evaluation protocol that distinguishes moment localization, answer correctness, and citation support.
\end{enumerate}
The contribution is the design and characterization of the system. We do not claim a new foundation model, state-of-the-art retrieval accuracy, or validated operation at archive scale.

\section{Related Work}
\label{sec:related}
\subsection{Temporal localization and video question answering}
QVHighlights provides free-form queries, annotated moments, and saliency labels, and introduces Moment-DETR, which predicts intervals directly \cite{qvhighlights}. FalconVQA instead retrieves precomputed chunks. This enables reuse across questions, but observed events can begin or end inside a returned chunk. NExT-QA emphasizes temporal and causal reasoning about actions \cite{nextqa}. EgoSchema further distinguishes video length from the amount of evidence needed to answer a question \cite{egoschema}. These settings motivate testing not only whether the answer is correct, but also whether the cited portion is sufficient to verify it.

\subsection{Retrieval-augmented generation for video}
Retrieval-augmented generation conditions a generator on documents from an external index \cite{lewis}. Video-RAG retrieves auxiliary text from audio, OCR, and object detections and feeds it to a video-language model alongside video frames \cite{videorag}. FalconVQA shares the use of derived evidence but targets repeated queries over an indexed repository. Its backend synthesis operates on stored descriptions, transcripts, and aggregates; the application returns the original footage separately for inspection. This avoids decoding it during retrieval, but limits the answer to information captured at ingest time.

\section{System Architecture}
\label{sec:architecture}
\begin{figure*}[!t]
\centering
\begin{tikzpicture}[
    node distance=8mm,
    cell/.style={draw=black!30, fill=blue!3, rounded corners=2pt,
      align=center, minimum width=30mm, minimum height=13mm, font=\small},
    store/.style={cell,fill=cyan!9},
    flow/.style={-{Stealth[length=2.2mm]},semithick,draw=black!70},
    link/.style={flow,dashed},
    lane/.style={font=\footnotesize\bfseries,anchor=east,xshift=-3mm}]
  \node[cell] (source) {\textbf{Video source}\\File or URL};
  \node[cell,right=of source] (grid) {\textbf{Ingest + segment}\\Shared timeline};
  \node[cell,right=of grid] (evidence) {\textbf{Evidence}\\Frame text + speech};
  \node[store,right=of evidence] (index) {\textbf{VIDRA artifacts}\\Index + aggregates};
  \node[cell,below=of source] (prompt) {\textbf{User question}\\Project context};
  \node[cell,below=of grid] (tools) {\textbf{Agent tools}\\Resolve video scope};
  \node[cell,below=of evidence] (search) {\textbf{Query service}\\Rank or synthesize};
  \node[cell,below=of index] (play) {\textbf{Workspace}\\Answer + clip reel};
  \node[lane,left=of source,node distance=0pt] {OFFLINE};
  \node[lane,left=of prompt,node distance=0pt] {ONLINE};
  \draw[flow] (source) -- (grid);
  \draw[flow] (grid) -- (evidence);
  \draw[flow] (evidence) -- (index);
  \draw[flow] (prompt) -- (tools);
  \draw[flow] (tools) -- (search);
  \draw[flow] (search) -- (play);
  \draw[link] (index.south west) -- (search.north east);
\end{tikzpicture}
\caption{FalconVQA separates ingestion from interactive search. The dashed
arrow denotes access to indexed artifacts. The workspace opens the original
media at the returned interval.}
\label{fig:architecture}
\end{figure*}


Fig.~\ref{fig:architecture} separates the ingestion and query paths. The application manages accounts, projects, conversation history, and per-project video records in Supabase. A FastAPI service accepts ingest jobs, exposes analytical artifacts, and serves moment search and question answering over VIDRA. Media and analytical artifacts are written to local storage, or mirrored to Azure Blob Storage when configured. The deployed backend uses local disk as a cache for remote analytical artifacts. The application database is not the active vector index.

\subsection{Ingestion, identity, and lifecycle}
A source is either a direct video URL or an uploaded file. In the latter case the browser obtains a temporary upload URL, writes the bytes to storage, and submits the upload identifier to the backend. The service identifies media by a truncated SHA-256 digest of its bytes; this identifier names the analytical record. The application maintains a separate video row and uses the ingest job result to associate it with the analytical record. During processing, the application polls the backend and reconciles job progress or failure into the project view. A completed record provides the playback URL, duration, and availability of analyses.

The backend serializes ingest jobs in one worker. It persists job state but does not persist the executable queue: work interrupted by a process restart is not automatically resumed. The deployment separates media transfer from the application server and reuses derived artifacts for queries. It does not imply that ingestion scales horizontally or that first-use retrieval avoids storage latency.

\subsection{Project scope and playback}
Retrieval tools resolve the user's ready recordings in the active project before passing explicit video identifiers to the backend. A user can tag videos to guide the model's selection; the server-side project check, not the tag or prompt, establishes which project videos can be retrieved. The backend itself does not manage tenant identities. The application and database therefore enforce project ownership, while the backend supports a deployment-level service token. Direct media access is a separate control and should not be inferred from project isolation in the chat interface.

A retrieved moment carries a recording identifier and start and end times in media seconds. The clip reel seeks to the start of each interval on the source media and advances at its end. The application can replay several ranges from one source without pre-rendering each one as a separate file. This makes a retrieval result inspectable; it does not certify that a generated statement is correct.

\section{Methodology}
\label{sec:method}
\subsection{Segmentation and evidence extraction}
Let a recording $v$ have duration $T_v$ and a timeline with edges $0=b_0<b_1<\cdots<b_m=T_v$. Its $i$th chunk is $c_i=(v,[b_i,b_{i+1}))$, where $0\leq i<m$. The key $(v,i)$ is necessary because chunk indices are local to a recording. The current backend maps the user's ingest settings to one VIDRA segmentation policy per run (Table~\ref{tab:policies}). In weighted mode, the strongest requested signal selects a policy; the deployed adapter does not fuse the four signals used in the earlier backend. If the selected policy requires an unavailable media stream, ingestion can fall back to a fixed grid.

\begin{table}[!t]
\caption{Segmentation Policies in the Current Pipeline}
\label{tab:policies}
\centering
\small
\begin{tabularx}{\columnwidth}{@{}lX@{}}
\toprule
\textbf{Policy} & \textbf{Source of interior boundaries} \\
\midrule
Uniform & Arithmetic time grid; no content detector. \\
Scene & Frame-to-frame color and brightness changes in the sampled visual stream. \\
VAD & Midpoints of sufficiently long gaps between transcript speech spans. \\
Speaker & Changes in speaker-labeled turns when those turns are available. \\
\bottomrule
\end{tabularx}
\end{table}

Visual samplers keep frames at regular times or according to cues associated with people, objects, or visible text. The vision-language model produces short descriptions for the selected frames in each chunk. The audio branch aligns Whisper-derived text to the same grid \cite{whisper}; speaker-attributed analysis can additionally use pyannote \cite{pyannote}. If a modality is absent or disabled, the remaining branch can still produce evidence. Neither sampling nor transcription guarantees coverage of brief or ambiguous events, and an event split by a boundary can be harder to express in a single unit.

The pipeline stores a timeline, derived descriptions and transcript artifacts when available, and an embedding index. It then computes video-level products from the derived artifacts rather than replaying the source. Depending on available evidence and the configured computation tier, these can include coverage statistics, a summary, chapters, events, or recurring entities. These outputs support different user requests: ``When is the box moved?'' needs a localized span, while ``What are the main topics?'' benefits from an overview.

\subsection{Dense unit ranking and moment aggregation}
In the current configuration the text representations are embedded in a shared space. The library supports optional visual embeddings and other storage implementations; the deployed adapter uses the dense-only, file-backed \texttt{Folder} implementation without a visual embedding role. For a query embedding $\mathbf{q}$ and a nonzero unit embedding $\mathbf{x}_j$, the backend ranks units by cosine similarity:
\begin{equation}
\operatorname{sim}(\mathbf{q},\mathbf{x}_j)=
\frac{\mathbf{q}^{\mathsf T}\mathbf{x}_j}
{\lVert\mathbf{q}\rVert_2\,\lVert\mathbf{x}_j\rVert_2}.
\label{eq:cosine}
\end{equation}
The candidate pool is truncated before hits are grouped by recording and chunk. Let $r_1(c)$ and $r_2(c)$ denote the one-based global ranks of a chunk's two highest-ranked units among the returned candidates. Inspired by reciprocal rank fusion \cite{rrf}, the implementation scores a moment as
\begin{equation}
S(c)=\frac{1}{10+r_1(c)}+
\begin{cases}
\dfrac{0.5}{10+r_2(c)} & \text{if a second hit exists},\\[3pt]
0 & \text{otherwise}.
\end{cases}
\label{eq:rank}
\end{equation}
Limiting the contribution to two hits prevents a chunk from accumulating score from every stored description. The second hit is not required to come from a different modality. The resulting score depends on the candidate ranking and is not a relevance probability or an absolute similarity threshold. In particular, a rank-one moment exists even for an unrelated query when the index supplies candidates.

\subsection{Synthesis and source presentation}
The search path returns ranked moments without requiring an answer. The question-answering path supplies full moment evidence to a text model and also consults available summaries and chapters for up to five scoped videos. The model is instructed to restrict its answer to the evidence, attach minute-and-second citations, and say when evidence is insufficient. The chat agent can also read a transcript, an individual chunk, or a specific aggregate without asking the backend to synthesize an answer every time. The application can show both the text response and a clip reel, with citations referring to timestamps on the original media timeline (Fig.~\ref{fig:evidence}).

\begin{figure}[!t]
\centering
\begin{tikzpicture}[
    x=1.08cm,y=0.76cm,
    smalltext/.style={font=\footnotesize},
    match/.style={fill=cyan!13,draw=blue!45,rounded corners=1pt}]
  \node[smalltext,anchor=west] at (0,1.55) {\textbf{Query:} When is the box moved?};
  \node[smalltext,anchor=west] at (0,0.94) {\textbf{Timeline of the returned recording}};
  \draw[gray!60] (0,0) rectangle (2,0.48);
  \draw[match] (2,0) rectangle (4,0.48);
  \draw[gray!60] (4,0) rectangle (6,0.48);
  \node[smalltext] at (3,0.24) {Matched chunk};
  \foreach \x/\t in {0/0:00,2/0:20,4/0:40,6/1:00}
    \node[smalltext] at (\x,-0.24) {\t};
  \node[smalltext,anchor=west] at (0,-0.88) {\textbf{Clip reel:} Seek to 0:20 and play through 0:40.};
\end{tikzpicture}
\caption{Illustrative citation-to-playback mapping. The interval is a returned
chunk, not a prediction of the exact onset and end of the action.}
\label{fig:evidence}
\end{figure}

A citation is a useful link, not an entailment check. The model can still cite an irrelevant moment, mistake a scene description for an observation, or draw an inference not supported by the footage. For example, footage may support that a box was moved without revealing who owned it. We therefore treat answer correctness, citation support, and abstention as distinct outcomes in the proposed evaluation.

\section{Evaluation}
\label{sec:evaluation}
We report two distinct evidence sets: a reproducible component pilot for the current VIDRA backend and historical results from the earlier prototype. They differ in both method and scope and are not successive rows in a controlled comparison. We close with the protocol for an end-to-end evaluation of the current system.

\subsection{Current backend: component pilot}
\label{sec:pilot}
The pilot used VIDRA 3.0.0 at repository revision \texttt{39081b9}, an Apple M2 with 8~GB of memory, macOS 26.0.1, and Python 3.12.14. The script generates its own footage and index. Neither experiment calls hosted models, uses private footage, or requires a production service. Reported p95 values are empirical percentiles of the timed runs, not bounds on deployed latency.

\paragraph{Boundary placement.}
We encoded three silent, 60-s H.264 videos at 320$\times$180 and 20 frames/s. Each alternates between dark and light frames at three known, deliberately off-grid transition times. The transitions align with the default scene detector's 250-ms sampling grid. We compared the default fixed 20-s grid with the default scene policy. A predicted interior boundary is a match when it is within 250~ms of an unmatched reference boundary. With $M$ matches, $P$ predictions and $G$ references, pooled boundary F1 is $2M/(P+G)$. We excluded media generation, media probing, and one warm-up per clip and policy; 30 observations per policy cover boundary evidence extraction and timeline generation.

\begin{table}[!t]
\caption{Current VIDRA: Boundary Pilot on Three Synthetic Videos}
\label{tab:boundaries}
\centering
\small
\begin{tabular}{@{}lrrrr@{}}
\toprule
\textbf{Policy} & \textbf{Hits/pred.} & $\boldsymbol{F_1}$ & \textbf{Median} & \textbf{p95} \\
\midrule
Uniform & 0/6 & 0.00 & 0.36~ms & 0.52~ms \\
Scene & 9/9 & 1.00 & 47.80~ms & 153.16~ms \\
\bottomrule
\end{tabular}
\end{table}

Table~\ref{tab:boundaries} shows that the scene policy found all nine synthetic cuts, while the fixed grid found none by construction. The scene policy has a higher cost because it decodes and compares frames; the uniform policy computes intervals without scoring the picture. Repeated timings increase the sample for latency, not the number of unique cut annotations. This trial is a positive control for the boundary detector, not evidence about event localization inside scenes.

\paragraph{File-backed search.}
We timed the current file-backed search and moment aggregation using 100, 1,000, and 5,000 synthetic units, with one seeded, 1,536-dimensional vector per chunk and 100 chunks per synthetic recording. We provided query vectors directly and checked that a vector's own chunk ranked first. Each query reads the index, scores the vectors, and reduces up to 20 candidates to five moments. There were 15 timed queries per size after one warm-up. These times include JSON parsing and ranking, but exclude query embedding, answer synthesis, network calls, and cloud cache misses.

\begin{table}[!t]
\caption{Current VIDRA: Warm Local Retrieval Latency}
\label{tab:retrieval}
\centering
\small
\begin{tabular}{@{}rrrr@{}}
\toprule
\textbf{Units} & \textbf{Videos} & \textbf{Median (ms)} & \textbf{p95 (ms)} \\
\midrule
100 & 1 & 12.2 & 15.4 \\
1,000 & 10 & 144.4 & 283.1 \\
5,000 & 50 & 770.9 & 1,154.3 \\
\bottomrule
\end{tabular}
\end{table}

Median search time increases with scope across the three measured operating points (Table~\ref{tab:retrieval}). This is consistent with the implementation reading and scoring all in-scope units; the data do not separate lookup, parsing, scoring, and allocation costs. The exact-vector check verifies that the search and moment mapping can return a known match. It is not a semantic accuracy measurement, and one unit per chunk does not exercise the relevance benefit of the second term in \eqref{eq:rank}.

\subsection{Historical results: predecessor backend}
\label{sec:legacy}
Before the VIDRA integration, the prototype used a different boundary pipeline and local embeddings in an embedded Qdrant index. Its development report documents seven publicly sourced fixtures, mainly short fixed-camera recordings, with one dialogue clip \cite{legacy}. Table~\ref{tab:legacy} preserves the reported exploratory findings. The historical report does not provide per-query raw scores or sample counts for these ranges, so we do not interpret them as estimates of population accuracy or report confidence intervals. Its top-three rule is an acceptance criterion, not a measured Recall@3 value. We likewise do not report the prototype's latency budgets as measured runtimes.

\begin{table}[!t]
\caption{Predecessor Only: Exploratory Findings From the Project Report \cite{legacy}}
\label{tab:legacy}
\centering
\small
\begin{tabularx}{\columnwidth}{@{}p{0.28\columnwidth}X@{}}
\toprule
\textbf{Check} & \textbf{Historically reported observation} \\
\midrule
Presence scoring & Observed present-content scores $\geq 0.665$ and absent-content scores $\leq 0.524$. \\
Person linking & Co-visible, known-different people scored as high as $0.941$; known-same pairs scored $0.889$--$0.915$. A single similarity threshold did not separate those cases. \\
Aggregate routing & Question-level mean scores ranged from $0.47$ to $0.56$; the prototype used relative selection. \\
\bottomrule
\end{tabularx}
\end{table}

The historical score ranges are not comparable to the current moment score in \eqref{eq:rank}. Their main value is to identify hypotheses to retest: whether absent-content queries get unwarranted matches, whether repeated observations are distinguishable from look-alikes, and whether summaries or chapters are helpful for broad questions.

\subsection{Protocol for real-video evaluation}
\label{sec:protocol}
\paragraph{Corpus and annotation.}
Use a public, license-compatible subset of QVHighlights for query-to-moment evaluation \cite{qvhighlights}. Evaluate single-video localization on its supplied scope, and add an explicitly documented repository setting with distractor recordings. To measure open-ended grounded answering, curate a separate collection that includes interviews or lectures, silent fixed-camera sequences, and videos requiring visual and speech evidence together. Split by source video, keeping calibration recordings separate from the test set. For each question, independent annotators mark answerability, an accepted answer or key claims, and one or more supporting intervals. Resolve disagreements before scoring. Include unanswerable questions and closely related distractors; publish counts, splits, and annotation instructions.

\paragraph{Comparisons and metrics.}
Compare transcript-only, vision-only, and combined evidence on matched recordings and settings. Compare fixed and scene segmentation where video is available, reporting chunk counts as well as accuracy. For QA, compare moments-only with moments plus summaries and chapters using the same generator and budget. A retrieval counts as a hit only if it is in the correct recording and its best temporal intersection over union (tIoU) with an annotated interval meets a prespecified threshold:
\begin{equation}
\operatorname{tIoU}(I,G)=
\frac{\operatorname{duration}(I\cap G)}{\operatorname{duration}(I\cup G)}.
\end{equation}
Report Recall@1 and Recall@5 at tIoU thresholds of 0.3 and 0.5, plus median absolute onset error for matched intervals. For generated answers, separately report answer correctness, the fraction of citations that support their associated claims, and the fraction of factual claims covered by a supporting citation. Report abstention on unanswerable questions and false abstention on answerable ones. Break down results by evidence modality and report paired intervals of uncertainty, resampling at the recording level.

\paragraph{Systems measurements.}
Log hardware, storage configuration, model identifiers and versions, corpus duration, and number of indexed units. Measure cold and warm query latency, time to playable evidence, ingest wall time per video minute, and provider call counts separately. Report median and p95 with the repetition count and concurrency. This protocol can distinguish a better retriever from an engine that is simply serving a smaller corpus or using a warm cache.

\subsection{Reproducibility}
The artifact accompanying this draft supplies the pilot script and its raw timing samples. To compare a later run, use the source revision and runtime versions recorded with the samples, and preserve both the benchmark configuration and observed output. The semantic-quality and end-to-end measurements described in Section~\ref{sec:protocol} were not run for this draft.


\section{Discussion}
\label{sec:discussion}
\subsection{What the measurements establish}
The current pilot shows that scene-derived boundaries can recover clean transitions missed by a fixed grid, at the cost of decoding and scoring frames. The result should not be extrapolated to gradual changes, handheld video, or actions inside an otherwise static shot. The local retrieval experiment shows that reading the file-backed index and ranking its vectors can itself become material as the search scope grows. It does not isolate file parsing from cosine computation or justify a specific corpus-size limit. A persistent in-memory representation or server-side index is a candidate comparison, not a measured improvement.

The earlier backend measurements describe engineering decisions made for a different index, model configuration, and score scale. They provide provenance for the project but cannot substitute for a matched comparison. In particular, transferring a similarity threshold from that backend to the rank-based score in \eqref{eq:rank} would be invalid. A head-to-head comparison should hold the footage, questions, evidence labels, and compute conditions fixed and report the two configurations explicitly.

\subsection{Threats to validity and deployment limits}
There are three main threats to validity. First, the pilot deliberately isolates components and uses synthetic inputs; it does not cover real visual or linguistic variation. Second, the retrieval timing uses warm runs on one machine, with media and index files on local storage. Cloud cache misses, provider latency, concurrent demand, and memory pressure can change both median and tail behavior. Third, the previous prototype report is a secondary source for the historical figures; the current artifact does not reproduce its raw trials. The appropriate comparison is a new, pre-specified experiment, not an aggregate across the result sets.

The system can index sensitive speech and images even when video bytes are stored separately. Hosted vision calls may transmit selected frames, embedding calls may transmit extracted text, and answer generation may transmit retrieved passages. A deployment should verify that these transfers, storage retention, and media access are appropriate for the recordings. In settings with safety, legal, or privacy consequences, verify evidence from original footage and do not equate the retrieval of a similar event with an identification or finding.

\section{Conclusion}
\label{sec:conclusion}
FalconVQA indexes visual and spoken evidence on a common timeline and makes retrieved moments playable within a conversational repository interface. The current component pilot demonstrates accurate boundary recovery for deliberately clean transitions and quantifies the latency of local file-backed retrieval over three synthetic corpus sizes. It does not establish retrieval or answer quality on real footage. The accompanying evaluation protocol defines how to measure temporal localization, citation support, and abstention independently, so future updates can test whether the implemented workflow delivers verifiable evidence at repository scale.

\section*{Acknowledgment}
The authors thank Dr. Pratiba D and the Centre of Excellence in Computer Vision Research at RV College of Engineering for project guidance.

% Balance the references in this draft; adjust or remove after repagination.
\IEEEtriggeratref{6}
\begin{thebibliography}{00}

\bibitem{qvhighlights}
J. Lei, T. L. Berg, and M. Bansal, ``QVHighlights: Detecting moments and highlights in videos via natural language queries,'' in \emph{Adv. Neural Inf. Process. Syst.}, vol. 34, 2021, pp. 11846--11858.

\bibitem{nextqa}
J. Xiao, X. Shang, A. Yao, and T.-S. Chua, ``NExT-QA: Next phase of question-answering to explaining temporal actions,'' in \emph{Proc. IEEE/CVF Conf. Comput. Vis. Pattern Recognit.}, 2021, pp. 9777--9786.

\bibitem{egoschema}
K. Mangalam, R. Akshulakov, and J. Malik, ``EgoSchema: A diagnostic benchmark for very long-form video language understanding,'' in \emph{Adv. Neural Inf. Process. Syst.}, vol. 36, 2023.

\bibitem{lewis}
P. Lewis \emph{et al.}, ``Retrieval-augmented generation for knowledge-intensive NLP tasks,'' in \emph{Adv. Neural Inf. Process. Syst.}, vol. 33, 2020, pp. 9459--9474.

\bibitem{videorag}
Y. Luo \emph{et al.}, ``Video-RAG: Visually-aligned retrieval-augmented long video comprehension,'' in \emph{Adv. Neural Inf. Process. Syst.}, vol. 38, 2025.

\bibitem{whisper}
A. Radford \emph{et al.}, ``Robust speech recognition via large-scale weak supervision,'' in \emph{Proc. Int. Conf. Mach. Learn.}, vol. 202, 2023, pp. 28492--28518.

\bibitem{pyannote}
H. Bredin \emph{et al.}, ``pyannote.audio: Neural building blocks for speaker diarization,'' in \emph{Proc. IEEE Int. Conf. Acoust. Speech Signal Process.}, 2020, pp. 7124--7128.

\bibitem{rrf}
G. V. Cormack, C. L. A. Clarke, and S. B\"uttcher, ``Reciprocal rank fusion outperforms Condorcet and individual rank learning methods,'' in \emph{Proc. Int. ACM SIGIR Conf. Res. Develop. Inf. Retrieval}, 2009, pp. 758--759, doi: 10.1145/1571941.1572114.

\bibitem{legacy}
A. Anish and V. Verma, ``FalconVQA,'' project development report, \emph{Product Performance} and \emph{Dataset Creation} sections, repository file \texttt{latex/phases.tex}, rev. \texttt{76fc7ff}.

\end{thebibliography}

\EOD
\end{document}
